% !TeX TS-program = xelatex
% !Mode TeXUTF-8
%# -- codingutf-8 --
% !TeX encoding = UTF-8
%%%%%%%%%%%%%%%%%%%%%%%%%%%%%%%%%%%%%%%%%
% 答辩小组成员的名单：
\newcommand{\zuzhang}{答辩组长}
\newcommand{\zuyuanone}{余老一} % 如果组员人数
\newcommand{\zuyuantwo}{袁老二}
\newcommand{\zuyuanthree}{张三三}
\newcommand{\zuyuanfour}{李四四}
\newcommand{\zuyuanfive}{王五五}
\newcommand{\zuyuansix}{刘六六}
\newcommand{\zuyuanseven}{齐七七}% 如果组员人数不足七位，就将“组员七”删除。上同
%----------------------------------------
% 答辩记录、评阅意见和论文评定成绩在下面填写
%--------论文评定成绩---------------------
\newcommand{\dpingfen}{83} % 具体的分数在左端的{}中录入
\newcommand{\jl}{
%------------答辩记录----------------------
\centering\fbox{答辩组成员的名单在文件{\tt dabianjilu.tex} 中指定的位置分别录入}

这里留出了八个人名的位置，如果组员人数不够八人，那就将“$\{\text{名字}\}$”中的``名字"删除，使之成为"$\{\}$"。

\centering{\fbox{答辩记录的内容在文件{\tt dabianjilu.tex} 中的这个位置录入}}
\begin{description}
	\item[问：] 什么是非线性递推数列？线性递推数列是怎么定义的？例子
	\item[答：]非线性递推数列就是递推关系为分式型或者含有根号和次方的数列。形如$a_{n+1} = \frac{aa_n+b}{ca_n+d}$
	类型的，线性递推数列就是含有$a_{n+1} = \uplambda_1a_1+\uplambda_2a_2+\cdots+\uplambda_na_n$ 递推关系式的数列。一般类型有$a_{n+1} = pa_n+q$。
	\item[问：] 有哪些形容词可以修饰递推数列？
	\item[答：] 高阶。（可以用齐次非齐次、常系数变系数、几阶）
\end{description}
\begin{center}
	\colorbox{yellow}{$\text{记录的字数不要超过预留空格所允许的数量}$}
\end{center}
%-------------------------------------------
}
\newcommand{\dcomments}{
%----------------------评审意见---------------------
\begin{center}
\fbox{对答辩的评审意见在文件{\tt dabianjilu.tex} 中的这个位置录入}
\end{center}		
\vspace{3mm}
	
本文结构合理，层次分明，思路清晰，显示出作者具有一定的综合运用所学知识的能力，符合本科毕业论文的要求，同意其毕业。
\begin{center}
	\colorbox{yellow}{$\text{评审意见的字数不要超出预留空格所允许的限额}$}
\end{center}
%----------------------------------------------------
}